\chapter{置換積分}\label{章:置換積分}

$x(x^2+1)^{20}$ を積分しようとすると，展開によって多くの項が生じる。しかし，式の中で繰り返し使われているまとまりは $x^2+1$ だけである。それを $t$ と書けば，べきの部分は $t^{20}$ になる。残った $x$ は邪魔なのだろうか。実は $x^2+1$ の微分は $2x$ なので，この $x$ も変数を換えるために必要な形になっている。

置換積分では，このように式の一部分を新しい変数にして，合成関数の微分を逆に使う。計算だけでなく，どの区間をどの区間に対応させたのかも確かめる。特に平方根や三角関数から元の変数を表すときには，一つの値に一つの値が対応する範囲を選ぶ必要がある。

\section{合成関数の微分を元へ戻す}\label{節:合成形の置換}

$F'(t)=f(t)$ であれば，合成関数の微分から
\[
 \frac{\dd}{\dd x}F(g(x))=f(g(x))g'(x)
\]
となる。この右辺の形が積分される式に現れたら，原始関数は左辺の $F(g(x))$ である。これが置換積分の根拠である。

ここで，微分可能で，その導関数も連続である関数を\textbf{連続微分可能}という。\label{用語:連続微分可能}

\begin{theorem}{合成形の置換積分}
区間 $I$ で $g$ は連続微分可能とし，$g$ の値を含む区間で $f$ は連続とする。$F'=f$ ならば
\[
 \int f(g(x))g'(x)\dd x=F(g(x))+C.
\]
さらに $a,b\in I$ に対して
\[
 \int_a^b f(g(x))g'(x)\dd x
 =\int_{g(a)}^{g(b)}f(t)\dd t.
\]
この二つの式に，$g$ が単調であるという条件は要らない。
\end{theorem}

不定積分の式は，上で確かめた合成関数の微分そのものである。定積分の式も，基本定理によって両辺がともに
\[
 F(g(b))-F(g(a))
\]
となることから分かる。$g$ が途中で増加から減少へ変わっても，この証明は変わらない。後で逆関数を使う置換と比べるため，この点を覚えておこう。

実際の計算では，$t=g(x)$ と置き，$\dd t=g'(x)\dd x$ と書く。この記号は，$g'(x)$ と $\dd x$ を含む積を，$t$ による積分へ書き換えることを表している。微小な数を自由に割り算するという説明に頼る必要はなく，根拠は上の合成関数の微分にある。

\begin{example}{まとまりとその微分を見つける}
$\int x(x^2+1)^{20}\dd x$ を求めよ。
\end{example}

$x^2+1$ を $t$ と置くと，$\dd t=2x\dd x$ だから，$x\dd x=\frac12\dd t$ とまとめて書き換えられる。したがって
\[
\begin{aligned}
 \int x(x^2+1)^{20}\dd x
 &=\frac12\int t^{20}\dd t\\
 &=\frac{t^{21}}{42}+C
 =\frac{(x^2+1)^{21}}{42}+C.
\end{aligned}
\]
新しい変数 $t$ で積分した後，不定積分の答えは元の変数 $x$ に戻す。右辺を微分すると
\[
 \frac{21}{42}(x^2+1)^{20}\,2x
 =x(x^2+1)^{20}
\]
となる。

ここでは $x=0$ でも式は成り立つ。計算の途中で $\dd x=\frac{\dd t}{2x}$ として $x$ で割る必要はなかった。$x\dd x$ 全体を $\frac12\dd t$ に換えれば，余分な除外点を作らずに済む。また $x^2+1$ は実数全体では単調でないが，原始関数の微分による確認には問題がない。

新しい変数を選ぶときは，「見た目が複雑な部分」だけでなく，「その部分を微分したものが残りにあるか」を調べる。たとえば
\[
 \int\frac{\log x}{x}\dd x
\]
では，$t=\log x$ と置けば $\dd t=\frac1x\dd x$ であり，$\int t\dd t$ に変わる。したがって，$x>0$ で答えは $\frac12(\log x)^2+C$ となる。複雑に見える対数を変数にすると，同時に分母の $x$ もなくなっている。

\section{分母の微分が分子にあるとき}\label{節:対数型の置換}

$\frac1t$ の原始関数は $\log|t|$ だから，置換積分を使うと
\[
 \int\frac{g'(x)}{g(x)}\dd x=\log|g(x)|+C
\]
となる。ただし，$g(x)\ne0$ である区間で使う。$g$ が連続ならば，その区間内で符号を変えるには $0$ を通る必要があるから，符号は一定である。正なら $\log g(x)$，負なら $\log(-g(x))$ を微分すれば，同じ商に戻る。

\begin{example}{絶対値を残すべき対数}
$\int\frac{6x}{3x^2-2}\dd x$ を求めよ。
\end{example}

分母の $3x^2-2$ を微分すると，分子の $6x$ が現れる。$t=3x^2-2$ と置けば $\dd t=6x\dd x$ なので，
\[
 \int\frac{6x}{3x^2-2}\dd x
 =\int\frac1t\dd t
 =\log|t|+C
 =\log|3x^2-2|+C.
\]
分母が $0$ になる $x=\pm\sqrt{\frac23}$ は除く。この二点の間では $3x^2-2<0$ なので，絶対値を外して $\log(3x^2-2)$ とすると実数の式ではなくなる。一方，$\log|3x^2-2|$ は三つの区間のどこでも使える形である。任意定数は各区間で独立に選んでよい。

分母が二次式だから必ず部分分数にする，という決め方はしない。この例では分子を見れば分母の微分があり，一度の置換で終わる。前章の式の整理と置換積分は競合する規則ではなく，式の特徴に応じて選ぶ道具である。

ただし，$\int\frac1{3x^2-2}\dd x$ の分子は $6x$ ではないから，同じ公式をそのまま使うことはできない。分母だけを $t$ と換えて，$\dd x$ を変えずに残すと，変数の異なる式を混ぜたことになる。置換後は，積分される式と微分の記号の両方を，新しい変数で書き直す必要がある。

\section{定積分では上下端も換える}\label{節:置換と積分区間}

不定積分では最後に $x$ へ戻したが，定積分なら新しい変数のまま計算を終えられる。$t=g(x)$ と置いたときは，$x=a$ に対応する $t=g(a)$ を下端に，$x=b$ に対応する $t=g(b)$ を上端にする。大きい数を自動的に上端にするのではなく，元の端点との対応を保つ。

たとえば $t=1-x$ で $x$ が $0$ から $1$ へ動けば，$t$ は $1$ から $0$ へ動く。さらに $\dd t=-\dd x$ だから，
\[
 \int_0^1(1-x)^2\dd x
 =-\int_1^0 t^2\dd t
 =\int_0^1t^2\dd t=\frac13.
\]
微分から来た負号と，上下端を入れ換えたときの負号は，それぞれ別の理由で現れる。途中で片方だけを落とさないよう，対応と微分を先に書いておくとよい。

\begin{example}{途中で折り返す置換}
$\int_{-1}^{2}2x\cos(x^2)\dd x$ を求めよ。
\end{example}

$t=x^2$ と置けば $\dd t=2x\dd x$ であり，端点の対応は
\[
 \begin{array}{c|cc}
 x&-1&2\\ \hline
 t&1&4
 \end{array}
\]
である。したがって合成形の置換積分から
\[
 \int_{-1}^{2}2x\cos(x^2)\dd x
 =\int_1^4\cos t\dd t
 =\sin4-\sin1.
\]
角 $1,4$ はラジアンである。$x$ が $-1$ から $2$ まで増える間に，$t=x^2$ は一度 $1$ から $0$ へ減り，その後 $4$ へ増える。端点だけを見た式が正しいのは，この折り返しが定積分の向きに反映されているためである。

図\ref{図:置換の折返し}に対応を示す。$x=0$ で積分を分けて確かめると，
\[
\begin{aligned}
 \int_{-1}^{2}2x\cos(x^2)\dd x
 &=\int_1^0\cos t\dd t+\int_0^4\cos t\dd t\\
 &=-\int_0^1\cos t\dd t+\int_0^4\cos t\dd t
 =\int_1^4\cos t\dd t.
\end{aligned}
\]
往復した範囲の寄与が符号付きで打ち消し合っている。この計算は，曲線の下の面積を重複も含めて足す操作とは違う。

\begin{figure}[htbp]
 \centering
 \begin{tikzpicture}[x=1.6cm,y=0.8cm]
  \draw[-{Stealth}] (-1.35,0)--(2.35,0) node[right] {$x$};
  \draw[-{Stealth}] (0,-0.35)--(0,4.65) node[above] {$t$};
  \draw[thick,blue,domain=-1:2,samples=65] plot (\x,{\x*\x});
  \draw[dashed] (-1,0)--(-1,1)--(0,1);
  \draw[dashed] (2,0)--(2,4)--(0,4);
  \fill (-1,1) circle (1.5pt);
  \fill (2,4) circle (1.5pt);
  \node[below] at (-1,0) {$-1$};
  \node[below right] at (0,0) {$0$};
  \node[below] at (2,0) {$2$};
  \node[left] at (0,1) {$1$};
  \node[left] at (0,4) {$4$};
  \node[right] at (1.6,1.9) {$t=x^2$};
  \draw[-{Stealth},thick] (-0.88,0.85)--(-0.55,0.41);
  \draw[-{Stealth},thick] (1.12,1.38)--(1.43,2.17);
 \end{tikzpicture}
 \caption{$x$ は増え続けるが，$t=x^2$ は減少してから増加する。}
 \label{図:置換の折返し}
\end{figure}

\section{元の変数を新しい変数で表す}\label{節:逆向きの置換}

合成形がそのまま見えない場合でも，根号などを新しい変数にしてから，元の変数を解き直すと式が簡単になることがある。このときは，$t=g(x)$ から $x$ を $t$ の関数として表せる範囲を選ぶ。

\begin{example}{根号を新しい変数にする}
$x>0$ において，$\int\frac{\sqrt{x}}{1+\sqrt{x}}\dd x$ を求めよ。
\end{example}

$t=\sqrt{x}$ と置けば，$t>0$ で $x=t^2$，$\dd x=2t\dd t$ となる。分子だけでなく分母と $\dd x$ も書き換えて
\[
 \int\frac{\sqrt{x}}{1+\sqrt{x}}\dd x
 =\int\frac{t}{1+t}\,2t\dd t
 =2\int\frac{t^2}{t+1}\dd t.
\]
根号はなくなったが分数が残った。ここで前章の割り算を使い，$t^2=(t+1)(t-1)+1$ と分ける。
\[
\begin{aligned}
 2\int\frac{t^2}{t+1}\dd t
 &=2\int\left(t-1+\frac1{t+1}\right)\dd t\\
 &=t^2-2t+2\log(t+1)+C.
\end{aligned}
\]
$t+1>0$ なので，ここでは対数の絶対値を外せる。$t=\sqrt{x}$ を戻すと，答えは
\[
 x-2\sqrt{x}+2\log(1+\sqrt{x})+C.
\]
微分して確かめると
\[
 1-\frac1{\sqrt{x}}+\frac1{(1+\sqrt{x})\sqrt{x}}
 =1-\frac1{1+\sqrt{x}}
 =\frac{\sqrt{x}}{1+\sqrt{x}}
\]
となる。置換が積分を一度で終わらせなくても，根号を有理式へ変えるという目的は達している。

この例で $x=t^2$ と書けたのは，$t=\sqrt{x}>0$ と決めていたからである。逆に $t=x^2$ から $x$ を求めるなら，$x=\sqrt{t}$ と $x=-\sqrt{t}$ の二つがある。$x>0$ の区間では前者，$x<0$ の区間では後者を使う。このように，逆関数を一つに決めるために選ぶ範囲を\textbf{枝}と呼ぶことがある。

区間上の連続関数で，逆関数を持つことと狭義単調であることは同値である。ここで狭義単調とは，異なる入力に対して出力が必ず一方向に増えるか，必ず一方向に減ることをいう。したがって，$t=g(x)$ から一つの関数 $x=h(t)$ を使って戻したいなら，そのような区間を選ばなければならない。さらに逆関数の微分を $h'(t)=\frac1{g'(h(t))}$ と書く場合には，分母が $0$ でないことも必要である。

これに対し，\sectionref{節:合成形の置換}の公式では逆関数を作っていない。$F(g(x))$ をそのまま微分したので，単調性を仮定しなかった。定積分の公式
\[
 \int_\alpha^\beta f(h(t))h'(t)\dd t
 =\int_{h(\alpha)}^{h(\beta)}f(x)\dd x
\]
自体も，同じ合成形の公式であり，$h$ の単調性を必要としない。単調性が必要になるのは，逆関数を一つ選んで式を書き直す段階である。この二つを区別すれば，置換に単調性が要る場合と要らない場合を混同せずに済む。

\section{三角関数で根号を簡単にする}\label{節:三角置換の初歩}

$a>0$ として，$\sqrt{a^2-x^2}$ には平方の差がある。三角関数の関係式 $1-\sin^2\theta=\cos^2\theta$ と形が同じだから，$x=a\sin\theta$ と置くことを考える。すると
\[
 \sqrt{a^2-x^2}
 =a\sqrt{1-\sin^2\theta}
 =a\sqrt{\cos^2\theta}
 =a|\cos\theta|.
\]
最後は $a\cos\theta$ ではなく，まず絶対値付きである。$\sqrt{z^2}=|z|$ だからである。そのうえで，たとえば $-\frac\pi2\le\theta\le\frac\pi2$ を選べば $\cos\theta\ge0$ となり，絶対値を外せる。この区間では $\sin\theta$ が一対一でもあるので，$-a\le x\le a$ の各 $x$ に対して $\theta$ が一つに決まる。

\begin{figure}[htbp]
 \centering
 \begin{tikzpicture}[x=2cm,y=1.5cm]
  \draw[-{Stealth}] (-1.95,0)--(1.95,0) node[right] {$\theta$};
  \draw[-{Stealth}] (0,-1.3)--(0,1.35) node[above] {$\sin\theta$};
  \draw[thick,blue,domain=-1.570796:1.570796,samples=70]
    plot (\x,{sin(\x r)});
  \draw[dashed] (-1.570796,0)--(-1.570796,-1)--(0,-1);
  \draw[dashed] (1.570796,0)--(1.570796,1)--(0,1);
  \node[above] at (-1.570796,0) {$-\frac\pi2$};
  \node[below] at (1.570796,0) {$\frac\pi2$};
  \node[above left] at (0,0) {$0$};
  \node[left] at (0,1) {$1$};
  \node[right] at (0,-1) {$-1$};
 \end{tikzpicture}
 \caption{正弦はこの区間で狭義単調増加し，逆向きの対応を一つに決められる。}
 \label{図:正弦の枝}
\end{figure}

\begin{example}{半円の根号を三角関数へ換える}
$a>0$ とする。$\int_0^a\sqrt{a^2-x^2}\dd x$ を求めよ。
\end{example}

ここでは $x\ge0$ だから，$x=a\sin\theta$ と置き，$0\le\theta\le\frac\pi2$ を選ぶ。$x=0$ は $\theta=0$，$x=a$ は $\theta=\frac\pi2$ に対応する。さらに
\[
 \dd x=a\cos\theta\dd\theta,\qquad
 \sqrt{a^2-x^2}=a\cos\theta
\]
なので，
\[
\begin{aligned}
 \int_0^a\sqrt{a^2-x^2}\dd x
 &=a^2\int_0^{\frac\pi2}\cos^2\theta\dd\theta\\
 &=a^2\left[\frac\theta2+\frac14\sin(2\theta)\right]_0^{\frac\pi2}
 =\frac{\pi a^2}{4}.
\end{aligned}
\]
曲線は $x^2+y^2=a^2$ の上半分であり，積分は半径 $a$ の円の四分の一の面積だから，幾何学的な結果とも一致する。積分計算そのものは，三角置換と二倍角の公式で得られている。

$0\le\theta\le\pi$ を選ぶと，$\sin\theta$ は $0$ から $1$ へ増えた後，$0$ へ戻る。この区間全体を $0\le x\le a$ と一対一に対応させることはできず，後半では $\cos\theta$ も負になる。枝を指定することは，こうした重複と符号を確定するために必要なのである。

もう一つ，$1+x^2$ には $1+\tan^2\theta=\frac1{\cos^2\theta}$ が対応する。$x=\tan\theta$ と置くと，微分も $\dd x=\frac1{\cos^2\theta}\dd\theta$ になる。この一致により，分母が消える場合がある。

\begin{example}{分母と微分の係数が打ち消し合う}
$\int_0^1\frac1{1+x^2}\dd x$ を求めよ。
\end{example}

$x=\tan\theta$ と置き，$0\le\theta\le\frac\pi4$ を選ぶ。この区間で正接は狭義単調増加し，端点は $x=0,1$ に対応する。よって
\[
\begin{aligned}
 \int_0^1\frac1{1+x^2}\dd x
 &=\int_0^{\frac\pi4}
   \frac1{1+\tan^2\theta}\,
   \frac1{\cos^2\theta}\dd\theta\\
 &=\int_0^{\frac\pi4}1\dd\theta
 =\frac\pi4.
\end{aligned}
\]
不定積分として表すなら，最後に $\tan\theta=x$ を満たす $\theta$ を $x$ の関数として書くことになる。この逆関数の名前，枝，微分は\chapterref{章:置換の背景}で詳しく扱う。ここでは定積分の端点が既知の角に対応するため，逆関数の記号を導入しなくても計算を終えられる。

\section{置換の後にも式を整える}\label{節:置換後の整理}

置換を選ぶとき，「すべてがすぐ基本公式になるか」だけを条件にすると，有用な置換を見落とす。根号をなくす，指数の中身を一つにする，分母の二次式を一次式にするなど，まず何を簡単にしたいかを決める。その後で，前章の式の整理を使えばよい。

\begin{example}{残った因子を新しい変数で表す}
$\int\frac{x^3}{1+x^2}\dd x$ を求めよ。
\end{example}

分母の $1+x^2$ を $t$ と置くと，$\dd t=2x\dd x$ である。分子には $x^3$ があるが，$x^3\dd x=x^2(x\dd x)$ と分ければ，$x\dd x=\frac12\dd t$ を使える。残る $x^2$ も $t-1$ と書けるので，
\[
\begin{aligned}
 \int\frac{x^3}{1+x^2}\dd x
 &=\frac12\int\frac{t-1}{t}\dd t\\
 &=\frac12\int\left(1-\frac1t\right)\dd t
 =\frac12t-\frac12\log t+C.
\end{aligned}
\]
$t=1+x^2>0$ を戻し，定数 $\frac12$ を任意定数へまとめると
\[
 \int\frac{x^3}{1+x^2}\dd x
 =\frac{x^2}{2}-\frac12\log(1+x^2)+C.
\]
$t=x^2+1$ は全実数上で一対一ではないが，残った $x^2$ を $t-1$ と書くのに逆関数 $x=h(t)$ は使っていない。$x$ の符号を選ぶ必要もない。

別の道もある。最初に割り算をすれば
\[
 \frac{x^3}{1+x^2}=x-\frac{x}{1+x^2}
\]
なので，第一項をそのまま積分し，第二項だけに置換を使って同じ答えを得る。どちらの場合も，式のどこを簡単にするかが分かっていれば，計算の途中で方針を見失わない。

これまでの置換では，複雑なまとまりとその微分が結び付いていた。しかし $x\e^x$ のように，二つの因子を片方の合成関数の微分としてまとめにくい積もある。この場合には，積の微分を逆に読む方法が役立つ。次章では，積分する因子と微分する因子を分けて考える。

\section*{章末問題}
\begin{enumerate}[label=\textbf{\thechapter.\arabic*},leftmargin=*]
 \item\label{問:置換べき} $\int(6x-4)(3x^2-4x+7)^5\dd x$ を求めよ。
 \item\label{問:置換二重対数} $\int\frac1{x\log x}\dd x$ を求めよ。実数の範囲で，どの区間に答えを適用できるかも述べよ。
 \item\label{問:置換非単調} $\int_{-2}^1 2x\e^{x^2}\dd x$ を求めよ。置換が単調でなくても成り立つ理由を説明せよ。
 \item\label{問:置換根号分母} $\int_0^3\frac1{1+\sqrt{1+x}}\dd x$ を求めよ。
 \item\label{問:置換円弧} $a>0$ とする。$x=a\sin\theta$ と置いて，$\int_0^{\frac a2}\sqrt{a^2-x^2}\dd x$ を求めよ。$\theta$ の区間も指定せよ。
 \item\label{問:置換正接} $x=\tan\theta$ と置いて，$\int_0^{\sqrt3}\frac1{1+x^2}\dd x$ を求めよ。
 \item\label{問:置換枝の違い} $t=x^2$ と置くと $x=-1,1$ がともに $t=1$ へ移る。このことから $\int_{-1}^1 1\dd x=0$ としてはいけない理由を述べよ。一方，$\int_{-1}^1\frac{2x}{1+x^2}\dd x=0$ は，同じ置換で正しく求められることを示せ。
\end{enumerate}

\smallskip\noindent 解答は\pageref{解答:第5章}ページ。
